
\begin{tikzpicture}[
    >=Stealth,
    sensor/.style={circle, draw=black!85, fill=gray!25, line width=1pt, minimum size=0.52cm},
    source/.style={circle, draw=black!85, fill=orange!75, line width=1.1pt, minimum size=0.52cm},
    axis/.style={->, line width=0.7pt, draw=black!60},
    velarrow/.style={->, line width=1.2pt, draw=orange!90!black}
]

% Eje unidimensional de movimiento
\draw[axis] (-5.6, 0.0) -- (5.6, 0.0)
    node[right=2pt, font=\small] {$x$ (trayectoria 1D)};

% Sensor R1
\coordinate (R1) at (-4.8, 0.0);
\node[sensor] (NodeR1) at (R1) {};
\node[below=7pt, align=center] at (R1) {
    \small\bfseries $R_1$\\[2pt]
    \footnotesize\color{black!65} $x=0$
};

% Sensor R2
\coordinate (R2) at (4.8, 0.0);
\node[sensor] (NodeR2) at (R2) {};
\node[below=7pt, align=center] at (R2) {
    \small\bfseries $R_2$\\[2pt]
    \footnotesize\color{black!65} $x=L$
};

% Fuente acústica móvil S(t)
\coordinate (S) at (0.0, 0.0);
\node[source] (NodeS) at (S) {};
\node[below=7pt, font=\small\bfseries] at (S) {$S(t)$};

% ---------------- FRENTES DE ONDA DOPPLER ----------------
% Hacia R1: longitudes de onda dilatadas
\draw[line width=0.8pt, red!75!black]
    ($(-0.1, 0.0)+(155:0.75)$) arc (155:205:0.75);
\draw[line width=0.8pt, red!75!black]
    ($(-0.35,0.0)+(155:1.45)$) arc (155:205:1.45);
\draw[line width=0.8pt, red!75!black]
    ($(-0.6, 0.0)+(155:2.15)$) arc (155:205:2.15);

% Hacia R2: longitudes de onda comprimidas
\draw[line width=0.8pt, blue!75!black]
    ($(-0.1, 0.0)+(-22:0.55)$) arc (-22:22:0.55);
\draw[line width=0.8pt, blue!75!black]
    ($(-0.35,0.0)+(-22:1.05)$) arc (-22:22:1.05);
\draw[line width=0.8pt, blue!75!black]
    ($(-0.6,0.0)+(-22:1.55)$) arc (-22:22:1.55);

% ---------------- VECTOR VELOCIDAD ----------------
\draw[velarrow] (0.0, 0.85) -- ++(1.15, 0.0)
    node[above=1pt, font=\small\bfseries] {$v(t)$};

\draw[dotted, line width=0.5pt, draw=orange!80!black]
    (0.0, 0.25) -- (0.0, 0.85);

% ---------------- COTAS DE LONGITUD DE ONDA ----------------
\draw[dotted, line width=0.5pt, red!75!black]
    (-2.0, 0.0) -- (-2.0, -0.9);
\draw[dotted, line width=0.5pt, red!75!black]
    (-1.1, 0.0) -- (-1.1, -0.9);

\draw[<->, draw=red!75!black, line width=0.6pt]
    (-2.0, -0.9) --
    node[below=2pt, font=\footnotesize] {$\lambda_1>\lambda_0$}
    (-1.1, -0.9);

\draw[dotted, line width=0.5pt, blue!75!black]
    (0.55, 0.0) -- (0.55, -0.9);
\draw[dotted, line width=0.5pt, blue!75!black]
    (1.05, 0.0) -- (1.05, -0.9);

\draw[<->, draw=blue!75!black, line width=0.6pt]
    (0.55, -0.9) --
    node[below=2pt, font=\footnotesize] {$\lambda_2<\lambda_0$}
    (1.05, -0.9);

% ---------------- RECUADRO DOPPLER ----------------
\node[
    draw=black!70,
    line width=0.7pt,
    fill=white,
    rounded corners=3pt,
    inner sep=6pt,
    anchor=south
] at (0.0, 2.05) {
    \footnotesize
    \begin{tabular}{@{} l l @{}}
        \multicolumn{2}{@{} l @{}}{
            \textbf{Efecto Doppler para fuente móvil ($v\ll c$):}
        } \\[3pt]

        Frecuencia en $R_1$ (alejamiento):
        & $f_1(t)\approx f_0\left(1-\frac{v(t)}{c(T)}\right)$ \\[2pt]

        Frecuencia en $R_2$ (acercamiento):
        & $f_2(t)\approx f_0\left(1+\frac{v(t)}{c(T)}\right)$ \\[4pt]

        \multicolumn{2}{@{} l @{}}{
            \textcolor{UdeAGreen!85!black}{
                \textbf{Estimador diferencial de velocidad:}
            }
            \quad
            $v(t)=c(T)\dfrac{f_2(t)-f_1(t)}{f_1(t)+f_2(t)}$
        } \\[2pt]

        \multicolumn{2}{@{} l @{}}{
            \textit{\scriptsize La frecuencia emitida se cancela algebraicamente en este modelo.}
        }
    \end{tabular}
};

\end{tikzpicture}
